\chapter{MARCO TEÓRICO}

\addcontentsline{toc}{section}{Introducción al marco teórico}

Para entender cómo funciona BONAIS hay que recorrer un camino, y ese
camino es el que traza este capítulo. Se empieza por lo más básico:
qué es una onda sonora y cómo viaja por el aire. Después se pregunta
qué ocurre cuando no una, sino dos ondas viajan al tiempo por el mismo
espacio, y se llega al fenómeno de la interferencia, con sus zonas de
refuerzo y de cancelación. Ese fenómeno se expresa después en la
escala con la que realmente se mide el sonido, el decibel, y en una
advertencia que se vuelve central para todo lo demás: los niveles no
se suman como números comunes.\\

\noindent Con esa física ya planteada, el capítulo da un giro hacia las
matemáticas: se muestra que el campo sonoro de un arreglo de fuentes
no es más que una función de dos variables, del tipo que estudia el
Cálculo III, y que las curvas de nivel, el gradiente y los puntos
críticos de esa función son, ni más ni menos, las herramientas que
permiten leer el mapa de sonido que genera BONAIS. Por último, el
capítulo se detiene en la industria: en que este problema ya lo
enfrentan a diario los ingenieros de sonido en vivo, y en cómo lo han
resuelto con los llamados arreglos direccionales de subwoofers, la
aplicación concreta que este proyecto busca modelar y optimizar.

\section{Cómo se propaga el sonido: la ecuación de onda}

La acústica estudia la generación, propagación y recepción de energía
mecánica en forma de ondas de presión que viajan a través de un medio
elástico \cite{kinsler2000,rossing2014handbook}. En un fluido como el
aire, las ondas sonoras son longitudinales: las partículas del medio
oscilan en la misma dirección en que se propaga la perturbación,
generando regiones alternadas de compresión y rarefacción
\cite{kinsler2000}.\\

\noindent Al aplicar la conservación de masa (ecuación de continuidad), la
segunda ley de Newton (ecuación de Euler) y la ecuación de estado a un
elemento de fluido, y linealizar el resultado suponiendo que las
perturbaciones acústicas son de amplitud pequeña frente a las
condiciones de equilibrio del medio, se obtiene la \emph{ecuación de
onda acústica lineal}, que distintos autores presentan de forma
equivalente en su versión tridimensional
\cite{kinsler2000,rossing2014handbook,pierce2014handbook}:

\begin{equation}
\nabla^2 p = \frac{1}{c^2}\frac{\partial^2 p}{\partial t^2},
\label{eq:onda3D}
\end{equation}

donde la presión acústica se representa matemáticamente como una función \(P(x,y,z,t)\) que depende de múltiples variables espaciales \((x,y,z)\) y del tiempo \(t\). La presencia del operador Laplaciano (\(\nabla^2\)) en la ecuación de onda tridimensional, el cual agrupa las segundas derivadas parciales espaciales, evidencia que la propagación del sonido es inherentemente un fenómeno multidimensional; esto justifica desde el primer momento la necesidad de aplicar el cálculo de varias variables para su análisis formal. Además, \(c=\sqrt{\mathcal{B}/\rho_0}\) es la
velocidad de propagación del sonido en el medio, con \(\mathcal{B}\) el
módulo de compresibilidad adiabático del fluido y \(\rho_0\) su densidad
de equilibrio; en el aire, a temperatura ambiente, $c\approx
343\ \text{m/s}$ \cite{kinsler2000,rossing2014handbook}.\\

Para señales
armónicas, con dependencia temporal \(e^{j\omega t}\), la
ecuación~\eqref{eq:onda3D} se transforma en la ecuación espacial de
Helmholtz \(\nabla^2 P + k^2 P = 0\), donde \(k=\omega/c=2\pi/\lambda\) es
el número de onda \cite{pierce2014handbook}.\\

La ecuación~\eqref{eq:onda3D} es la base de la que se derivan las
soluciones empleadas en este trabajo. Su solución más simple es la
onda plana armónica; sin embargo, para una fuente aislada y de tamaño
pequeño frente a la longitud de onda (el caso de un subwoofer, que
interesa a este proyecto), la solución relevante es la \emph{onda
esférica}. Cuando la fuente radia por igual en todas direcciones,
conviene expresar la ecuación~\eqref{eq:onda3D} en coordenadas
esféricas; para el caso de simetría radial, su solución armónica
divergente es

\begin{equation}
p(r,t) = \frac{A}{r}\, e^{j(\omega t - kr)},
\label{eq:onda_esferica}
\end{equation}

es decir, una onda cuya amplitud decae como \(1/r\) a medida que se aleja
de la fuente, a diferencia de la amplitud constante de la onda plana
\cite{kinsler2000}. Al introducir explícitamente una fuente puntual
ubicada en \(\vec{r}_0\), la presión que radia se escribe como

\begin{equation}
p(\vec{r},t) = \frac{A}{|\vec{r}-\vec{r}_0|}\,
\exp\!\big[j(\omega t - k|\vec{r}-\vec{r}_0|)\big].
\label{eq:fuente_puntual}
\end{equation}

¿Por qué es razonable tratar un subwoofer, que es un objeto físico con
dimensiones, como si fuera un simple punto matemático? La respuesta
está en comparar su tamaño con la longitud de onda que emite. En el
rango de interés del sonido en vivo para subwoofers
(\(40\ \text{Hz}\le f\le 100\ \text{Hz}\)), la longitud de onda
\(\lambda=c/f\) oscila entre \(3{,}4\ \text{m}\) y \(8{,}6\ \text{m}\). Frente
a esas dimensiones, las de la caja acústica de un subwoofer, del orden
de $0{,}6$ a \(0{,}8\ \text{m}\), son notablemente menores, de modo que
se cumple la condición de \emph{fuente acústicamente compacta},
\(kd_{\text{caja}}\ll 1\) \cite{pierce2014handbook}. Bajo este régimen,
la radiación carece de directividad propia a nivel de caja y el
subwoofer se comporta, en muy buena aproximación, como la fuente
puntual (monopolo esférico) de la
ecuación~\eqref{eq:fuente_puntual}, radiando de manera casi
omnidireccional \cite{pierce2014handbook,rossing2014handbook}. Esta
formulación es la que se combina linealmente, como se muestra a
continuación, para construir los campos de interferencia de un
arreglo de varias fuentes \cite{kinsler2000}.


\noindent La figura~\ref{fig:onda_esferica} muestra cómo se ve esta solución: frentes de onda circulares separados una longitud de onda, cuya amplitud disminuye a medida que se alejan de la fuente.

\begin{figure}[H]
\centering
\includegraphics[width=\textwidth]{figuras/mt_onda_esferica.png}
\caption{Izquierda: frentes de onda de una fuente puntual a 80~Hz. Derecha: el nivel cae 6~dB cada vez que se duplica la distancia, por el decaimiento $1/r$. Fuente: elaboración propia.}
\label{fig:onda_esferica}
\end{figure}

\section{Cuando dos fuentes se encuentran: superposición e interferencia}
\label{sec:interferencia}

La ecuación de onda~\eqref{eq:onda3D} es lineal, lo que implica que si
\(p_1\) y \(p_2\) son soluciones válidas (por ejemplo, los campos de
presión producidos por dos fuentes independientes), entonces su suma
\(p=p_1+p_2\) también lo es. A esta propiedad se le conoce como
\emph{principio de superposición} y constituye la base matemática de
todos los fenómenos de interferencia acústica
\cite{kinsler2000,rossing2014handbook}.\\

Un ejemplo representativo, desarrollado en detalle por
Kinsler et al.\ \cite{kinsler2000}, es el del \emph{doblete o dipolo
acústico}: dos fuentes puntuales idénticas, separadas una distancia
\(d\) pequeña frente a la longitud de onda (\(kd\ll 1\)), que vibran en
fase opuesta (desfasadas \(180^{\circ}\)). Sumando las contribuciones de
cada fuente, cada una de la forma de la
ecuación~\eqref{eq:fuente_puntual}, se obtiene en el campo lejano

\begin{equation}
p(r,\theta,t) = -j\,\frac{Akd}{r}\,\operatorname{sen}\theta\;
e^{j(\omega t - kr)},
\label{eq:dipolo}
\end{equation}

de modo que la amplitud sigue un patrón proporcional a
\(|\operatorname{sen}\theta|\): el campo se refuerza en el plano
perpendicular al eje de las fuentes y se cancela por completo a lo
largo de dicho eje \cite{kinsler2000}. Este ejemplo ilustra cómo la
interferencia entre dos ondas coherentes (misma frecuencia, diferencia
de fase constante) produce patrones espaciales de refuerzo y
cancelación que dependen de la geometría relativa de las fuentes y del
punto de observación. Cabe precisar que la
ecuación~\eqref{eq:dipolo} se dedujo bajo la aproximación \(kd\ll 1\); el
arreglo de subwoofers estudiado en este proyecto opera, en cambio, en
el régimen \(kd=\pi/2\) (separación de un cuarto de longitud de onda),
donde dicha aproximación \textbf{no} es válida. Por esta razón, el
modelo de BONAIS no usa la forma simplificada del dipolo, sino la
superposición exacta (sin aproximar) de las ondas esféricas de cada
fuente, aplicando únicamente el principio de superposición que este
ejemplo ilustra.\\


\noindent La figura~\ref{fig:interferencia} resume esta idea: en el tiempo, dos ondas en fase se suman y en contrafase se anulan; en el plano, eso se traduce en lóbulos de refuerzo y franjas de cancelación.

\begin{figure}[H]
\centering
\includegraphics[width=\textwidth]{figuras/mt_interferencia.png}
\caption{Izquierda: dos ondas en fase se refuerzan y en contrafase se cancelan. Derecha: patrón de refuerzo y cancelación de dos fuentes en fase separadas 6~m a 80~Hz. Fuente: elaboración propia.}
\label{fig:interferencia}
\end{figure}


Ese principio se generaliza sin dificultad a un arreglo de \(N\) fuentes,
cada una con su propia amplitud, y con la posibilidad de aplicarle a
cada una un retardo electrónico \(\tau_n\) y una inversión de polaridad,
ambos ajustables mediante un procesador de señal:

\begin{equation}
P_{\text{total}}(x,y) = \sum_{n=1}^{N} \frac{A_n}{r_n(x,y)}\,
e^{-j\left(k\,r_n(x,y) + \phi_n\right)},
\qquad
r_n(x,y) = \sqrt{(x-x_n)^2+(y-y_n)^2},
\label{eq:campo_total}
\end{equation}

donde \(\phi_n=\omega\tau_n+\theta_n\) reúne la fase introducida por el
retardo \(\tau_n\) y una eventual inversión de polaridad
\(\theta_n\in\{0,\pi\}\) \cite{pierce2014handbook}. La diferencia de fase
neta con la que dos fuentes llegan a un punto del espacio, y que
determina si allí hay refuerzo o cancelación, queda entonces
determinada tanto por la diferencia de caminos recorridos como por el
retardo y la polaridad aplicados electrónicamente:

\begin{equation}
\Delta\Phi(x,y) = k\big(r_2(x,y)-r_1(x,y)\big) + \omega\tau + \Delta\phi.
\label{eq:fase_neta}
\end{equation}

Esta última observación es la que permite el paso de la física a la
ingeniería: si la diferencia de caminos \(r_2-r_1\) ya está fijada por
la geometría del arreglo, el retardo \(\tau\) y la polaridad quedan como
\emph{parámetros de diseño} que un ingeniero puede ajustar a voluntad
para inducir interferencia constructiva hacia una dirección deseada
(el público) y destructiva hacia otra (el escenario), dando origen al
arreglo cardioide que se describe más adelante en este capítulo.\\

De la ecuación~\eqref{eq:campo_total} se obtiene además el campo de
presión eficaz al cuadrado, proporcional a la energía e independiente
del tiempo, \(p_{\text{rms}}^2(x,y)=\tfrac{1}{2}|P_{\text{total}}(x,y)|^2\)
\cite{pierce2014handbook}. Esta función bidimensional es, precisamente,
la base analítica sobre la que se construyen las curvas de nivel, el
gradiente y las integrales de energía que emplea BONAIS, y que se
presentan en las secciones siguientes.

\section{El decibel explica por qué los niveles no se suman}
\label{sec:decibel}

Dado que las magnitudes acústicas de interés (potencia, intensidad,
presión) abarcan un intervalo enorme de valores, es habitual
expresarlas en una escala relativa y logarítmica. Si \(W_1\) es la
potencia de una fuente y \(W_0\) una potencia de referencia, el nivel
correspondiente, en belios, se define como
\(L_1=\log_{10}(W_1/W_0)\); como el belio resulta una unidad
impráctica, se emplea en su lugar el \emph{decibel}, una décima parte
de un belio \cite[pág.~27]{everest2001}:

\begin{equation}
L_1 = 10 \log_{10}\!\left(\frac{W_1}{W_0}\right) \quad \text{[dB]}.
\label{eq:decibel_potencia}
\end{equation}

Esta definición aplica directamente a cualquier magnitud proporcional
a la energía o a la potencia, como la intensidad acústica \(I\). Puesto
que la intensidad de una onda es proporcional al cuadrado de la
presión acústica \(p\), al sustituir \(I\propto p^2\) en la
ecuación~\eqref{eq:decibel_potencia} el exponente al cuadrado se
convierte en un factor 2 delante del logaritmo, de modo que el nivel
expresado en términos de presión resulta \cite[pág.~27]{everest2001}:

\begin{equation}
L_p = 10 \log_{10}\!\left(\frac{p_1^2}{p_2^2}\right)
    = 20 \log_{10}\!\left(\frac{p_1}{p_2}\right) \quad \text{[dB]}.
\label{eq:decibel_presion}
\end{equation}

Para obtener un valor absoluto y no solo una comparación entre dos
presiones, se fija una presión de referencia común e
internacionalmente aceptada,
\(p_{\text{ref}}=0\text{,}0002\ \mu\text{bar}=20\ \mu\text{Pa}\), con la
cual la ecuación~\eqref{eq:decibel_presion} define el \emph{nivel de
presión sonora} (SPL) de cualquier sonido \cite[pág.~28]{everest2001}. Esta es
la magnitud que se reporta al medir un sonido con un sonómetro, y la
que se emplea en este proyecto en la forma equivalente
\(L=10\log_{10}(E/E_{\text{ref}})\), al trabajar con el campo de energía
\(E=|P|^2\) de un arreglo de fuentes (el mismo campo
\(p_{\text{rms}}^2(x,y)\) presentado al final de la sección anterior.\\

Una consecuencia directa de que el decibel sea una escala logarítmica
es que \textbf{los niveles no se pueden sumar aritméticamente}: el
nivel producido por dos fuentes sonando a la vez no es la suma de sus
niveles individuales en dB. Everest ilustra este punto con el caso de
dos fuentes de ruido de un estudio, un aire acondicionado con nivel de
\(55\ \text{dB}\) y un ventilador con nivel de \(60\ \text{dB}\): el nivel
combinado de ambos sonando al tiempo no es \(115\ \text{dB}\), sino que
se obtiene devolviendo cada nivel a su magnitud original (potencia o
intensidad relativa), sumando esas magnitudes, y solo entonces
tomando el logaritmo del resultado \cite[pág.~36]{everest2001}:

\begin{equation}
L_{\text{combinado}} = 10 \log_{10}\!\left(10^{L_1/10} + 10^{L_2/10}\right)
\quad \text{[dB]},
\label{eq:combinar_db}
\end{equation}

lo que en el ejemplo citado da $10\log_{10}(10^{5.5}+10^{6.0})\approx
61\text{,}2\ \text{dB}$, un valor apenas superior al mayor de los dos
niveles individuales y muy distinto de su suma aritmética
\cite[pág.~36]{everest2001}. La misma idea, aplicada en sentido inverso
(restando la contribución de una fuente conocida a un nivel total
medido), permite además aislar el nivel de una fuente que no se puede
medir por separado \cite[pág.~36]{everest2001}.\\

La ecuación~\eqref{eq:combinar_db} es, en el fondo, la versión en
decibeles del principio de superposición presentado en la sección
anterior: para combinar el efecto de varias fuentes hay que trabajar
primero con la magnitud física (presión, si las fuentes son coherentes
y hay que respetar su fase, como en la
ecuación~\eqref{eq:campo_total} y expresar el resultado en dB solo al
final. Esta observación es la que exige que, en todo el desarrollo
matemático que sigue, el campo total se calcule siempre como una suma
compleja de las contribuciones individuales antes de convertir
cualquier resultado a decibeles.

\section{El campo sonoro como función de dos variables}
\label{sec:campo_total}

\noindent Hasta aquí, el campo de presión de un arreglo de subwoofers,
\(P_{\text{total}}(x,y)\), se ha tratado como un objeto físico. Pero
basta con mirarlo de nuevo para notar algo importante: a cada punto
\((x,y)\) del plano (cada posición del público o del escenario) le
corresponde un único valor de energía o de nivel sonoro. Eso es,
exactamente, la definición de una función de dos variables.\\

Una función de dos variables es una regla de correspondencia que
asigna a cada par ordenado de números reales \((x,y)\) de un
subconjunto \(D\) del plano un único número real \(z=f(x,y)\)
\cite[cap.\ funciones de varias variables]{stewart2018},
\cite[cap.\ 12]{larson2015}. El dominio \(D\) representa el conjunto de
todas las posibles entradas, mientras que el rango es el conjunto de
todos los valores de salida \(z\). Geométricamente, la gráfica de una
función continua de dos variables representa una superficie
tridimensional en el espacio.\\

Para facilitar la interpretación de esa superficie sin necesidad de
dibujarla en tres dimensiones, se emplean las \emph{curvas de nivel}.
Una curva de nivel agrupa todos los puntos \((x,y)\) del dominio donde
la función toma un valor constante \(c\), es decir, todos los puntos que
satisfacen

\begin{equation}
f(x,y) = c.
\end{equation}

El conjunto de estas curvas, proyectadas sobre el plano \(xy\), conforma
un mapa de contorno que describe la topografía de la función: cuando
las curvas de nivel están muy juntas, la superficie es más empinada, y
cuando están separadas, la inclinación es menor.\\

En el diseño y análisis acústico de este proyecto, el campo sonoro se
representa exactamente como una función espacial de dos variables
\(L(x,y)\), la misma que se definió en la sección anterior a partir de
\(P_{\text{total}}(x,y)\). Este concepto matemático es el que permite
interpretar gráficamente el mapa de calor que genera BONAIS: las
líneas que en ese mapa separan zonas de distinto color son, por
definición, curvas de nivel de la función de presión sonora. A través
de esta representación se puede observar y predecir el comportamiento
espacial del campo, viendo con claridad cómo las zonas de refuerzo
(interferencia constructiva) y de cancelación (interferencia
destructiva) adquieren distribuciones geométricas medibles, que
típicamente forman curvas semejantes a hipérbolas. Esta base matemática
es la que sustenta de manera directa los objetivos específicos relacionados
con la evaluación de la energía acústica y las condiciones de separación
entre las fuentes.


\noindent La figura~\ref{fig:superficie_curvas} muestra las dos formas de ver la misma función: como superficie en el espacio y como mapa de curvas de nivel, que es la vista que usa BONAIS.

\begin{figure}[H]
\centering
\includegraphics[width=\textwidth]{figuras/mt_superficie_curvas.png}
\caption{El nivel $L(x,y)$ de un arreglo cardioide a 80~Hz visto como superficie (izquierda) y como curvas de nivel (derecha). Fuente: elaboración propia.}
\label{fig:superficie_curvas}
\end{figure}

\section{Gradiente, puntos críticos y el criterio del Hessiano}
\label{sec:gradiente}

Una vez que el campo sonoro se reconoce como una función de dos
variables, la siguiente pregunta natural es cómo cambia esa función de
un punto a otro, y dónde se encuentran sus valores extremos. Para
responderla, el Cálculo III ofrece un conjunto de herramientas que se
presentan a continuación en el mismo orden en que se necesitan.\\

Para estudiar la tasa de cambio de una función \(f(x,y)\) respecto a una
sola de sus variables se utilizan las \emph{derivadas parciales}
\cite{apostol,larson2015,stewart2018}. La derivada parcial respecto
a \(x\) se calcula derivando la función y tratando a \(y\) como una
constante, y de forma análoga para la derivada respecto a \(y\); se
denotan como \(f_x(x,y)\) o \(\partial f/\partial x\), y \(f_y(x,y)\) o
\(\partial f/\partial y\), respectivamente.\\

A partir de las primeras derivadas parciales se construye el vector
\emph{gradiente}, denotado \(\nabla f(x,y)\)
\cite[pág.~310]{apostol} \cite[pág.~1126]{larson2015}
. Este vector contiene las
derivadas parciales como sus componentes:

\begin{equation}
\nabla f(x,y) = \frac{\partial f}{\partial x}\,\hat{\imath}
              + \frac{\partial f}{\partial y}\,\hat{\jmath}.
\end{equation}

El gradiente tiene una propiedad geométrica fundamental, ya anticipada
al hablar de las curvas de nivel: siempre apunta en la dirección de
mayor tasa de incremento de la función, y es ortogonal a la curva de
nivel que pasa por cada punto \cite{marsden2011,stewart2018}.\\

En el marco de la optimización de funciones de varias variables, los \emph{puntos críticos} de la función analizada (como la presión acústica \(P\) o un campo genérico \(f\)) se hallan teóricamente en aquellas coordenadas del dominio donde el vector gradiente es exactamente igual a cero,

\begin{equation}
\nabla P = \vec{0},
\end{equation}

\noindent o donde las derivadas parciales no existen. Estos puntos críticos representan las ubicaciones candidatas a ser valores extremos de la superficie estudiada. Para determinar rigurosamente la naturaleza geométrica de ese punto crítico (es decir, si representa un máximo (refuerzo acústico), un mínimo (cancelación) o un punto de silla), la teoría exige el uso de las segundas derivadas parciales organizadas algebraicamente en la matriz Hessiana \cite[pág. 196]{marsden2011}, \cite[pág. 946]{stewart2018}, \cite[pág. 378 y 380]{apostol}. Sea \(D\) el discriminante
evaluado en un punto crítico \((a,b)\), definido a partir del
determinante del Hessiano:

\begin{equation}
D = f_{xx}(a,b)\,f_{yy}(a,b) - \big[f_{xy}(a,b)\big]^2.
\end{equation}

El criterio establece que:
\begin{itemize}
    \item si \(D>0\) y \(f_{xx}(a,b)>0\), entonces \(f(a,b)\) es un
    \textbf{mínimo local};
    \item si \(D>0\) y \(f_{xx}(a,b)<0\), entonces \(f(a,b)\) es un
    \textbf{máximo local};
    \item si \(D<0\), entonces \(f(a,b)\) es un \textbf{punto de silla}
    (no es máximo ni mínimo).
\end{itemize}

\noindent Este andamiaje teórico es, precisamente, el que fundamenta los
procesos internos que ejecuta el software BONAIS. Al analizar el
comportamiento de un arreglo, la aplicación necesita identificar los
picos y las caídas del nivel sonoro, permitiendo así mapear la
distribución de energía en el espacio de estudio, paso previo
indispensable para alcanzar los objetivos del proyecto. Para lograr un
procesamiento algorítmico numéricamente estable, la aplicación no deriva
directamente la función de nivel sonoro en dB (que, al ser logarítmica,
se dispara cerca de los nulos de presión), sino que aplica el gradiente
y el criterio del Hessiano sobre el campo de energía cuadrática \(E=|P|^2\)
definido en las secciones anteriores. De esta manera, cada vez que se explora
una coordenada en el mapa de BONAIS, la aplicación puede reportar de forma
eficiente la dirección de mayor cambio del campo y la localización
exacta de sus puntos críticos.


\noindent En la figura~\ref{fig:gradiente} cada flecha es el vector gradiente en ese punto: apunta hacia donde el nivel sube más rápido y corta perpendicularmente las curvas de nivel.

\begin{figure}[H]
\centering
\includegraphics[width=0.75\textwidth]{figuras/mt_gradiente.png}
\caption{Campo de gradiente del nivel para dos fuentes en fase separadas 5~m a 80~Hz. Cada flecha apunta hacia donde el nivel crece más rápido y es perpendicular a las curvas de nivel. Entre las fuentes hay un máximo local. Fuente: elaboración propia.}
\label{fig:gradiente}
\end{figure}

\subsection{Dos dominios distintos para una misma herramienta: el espacio y el diseño}

\noindent La sección anterior aplicó el gradiente y el criterio del Hessiano
sobre el campo \(E(x,y)\), es decir, sobre el \emph{dominio espacial}: a
cada coordenada del recinto le corresponde un valor de energía. Ese
análisis permite comprender el comportamiento del campo sonoro en el
espacio de estudio, pero \textbf{no} es, por sí solo, el que responde
al objetivo general del proyecto. Determinar la separación \(d\) y el
retardo \(\tau\) que maximizan la atenuación trasera exige aplicar
exactamente el mismo herramental (gradiente, puntos críticos, Hessiano)
pero sobre un \emph{dominio distinto}: ya no el de las coordenadas \((x,y)\)
del recinto, sino el de los \emph{parámetros de diseño} \((d,\tau)\) del arreglo.\\

\noindent Para ello se define una segunda función, la \emph{función de
atenuación}, que ya no depende de la posición sino únicamente de la
separación y el retardo elegidos para el arreglo. Sean
\(E_{\text{frente}}(d,\tau)\) y \(E_{\text{atrás}}(d,\tau)\) la energía del
campo, calculada mediante la ecuación~\eqref{eq:campo_total}, en un
punto de referencia situado respectivamente delante y detrás del
arreglo. Siguiendo la misma regla establecida en la
sección~\ref{sec:decibel} de trabajar siempre con la magnitud física
antes de tomar el logaritmo, la atenuación trasera del arreglo, en
decibeles, se define como

\begin{equation}
A(d,\tau) = 10\log_{10}\!\left(\frac{E_{\text{frente}}(d,\tau)}
{E_{\text{atrás}}(d,\tau)}\right) \quad \text{[dB]}.
\label{eq:funcion_atenuacion}
\end{equation}

\noindent La ecuación~\eqref{eq:funcion_atenuacion} es, formalmente, una función
de dos variables \(A=A(d,\tau)\), en todo comparable a la función
\(E(x,y)\) de la sección anterior. Por lo tanto, el objetivo específico
centrado en determinar la separación y el retardo que maximizan la
atenuación trasera se resuelve buscando los puntos \((d^\ast,\tau^\ast)\) donde

\begin{equation}
\nabla A(d^\ast,\tau^\ast) =
\left(\frac{\partial A}{\partial d},\ \frac{\partial A}{\partial \tau}\right)
= \vec{0},
\end{equation}

\noindent y clasificando dicho punto crítico con el mismo criterio del Hessiano
presentado en la sección anterior, ahora aplicado sobre
\(A(d,\tau)\): si el discriminante
\(D=A_{dd}A_{\tau\tau}-A_{d\tau}^2\) es positivo y \(A_{dd}<0\), el punto
\((d^\ast,\tau^\ast)\) corresponde a un \textbf{máximo} de la atenuación
trasera, que es precisamente lo que busca el objetivo general del
proyecto. El desarrollo explícito de \(A(d,\tau)\) y la determinación
numérica de \((d^\ast,\tau^\ast)\) se presentan en el capítulo del
modelo matemático, una vez definida la geometría concreta del arreglo.\\

\noindent Conviene subrayar la diferencia entre los dos análisis para no
confundirlos en el resto del documento: los puntos críticos de
\(E(x,y)\) (sección~\ref{sec:gradiente}) señalan \emph{dónde}, en el
recinto, el sonido se refuerza o se cancela, para un arreglo ya
construido; los puntos críticos de \(A(d,\tau)\) (esta sección) señalan
\emph{cómo construir} ese arreglo para que dicho refuerzo y
cancelación queden orientados de la manera más favorable posible.\\

\noindent Debe advertirse, además, que al igualar \(\nabla E(x,y)=\vec{0}\) o
\(\nabla A(d,\tau)=\vec{0}\) se obtiene, en general, un sistema de
ecuaciones no lineales que combina términos inversos \(1/r\) con sumas
de exponenciales complejas, para el cual no existe, salvo en casos muy
particulares, una solución analítica cerrada. Por esta razón, BONAIS
no resuelve estos sistemas a mano, sino mediante \textbf{métodos
numéricos de aproximación}: los puntos críticos espaciales se hallan
con un método iterativo tipo Newton, y la optimización de
\(A(d,\tau)\) se realiza mediante un método de descenso (o ascenso)
más pronunciado sobre el gradiente \cite{chongzak2013}. Esta es una
práctica estándar en optimización de funciones de varias variables
cuando la función objetivo no admite una solución cerrada simple.

\section{Integrales dobles sobre regiones generales}
\label{sec:integrales_dobles}

El primer objetivo específico de este proyecto busca establecer la
relación entre la energía total en una región del recinto (por
ejemplo, la zona ocupada por el público o la zona del escenario) y la
suma de las energías que cada fuente aportaría por separado. Para
plantear esa comparación con rigor es necesario definir, primero, qué
significa la energía \emph{total} en una región y no solo en un punto.\\

\noindent Dada una función continua \(f(x,y)\) definida sobre una región \(R\) del
plano, la integral doble

\begin{equation}
\iint_R f(x,y)\, dA
\end{equation}

\noindent se define, igual que la integral de una variable, como el límite de
una suma de Riemann: se cubre \(R\) con una malla de rectángulos
pequeños de área \(\Delta A\), se evalúa \(f\) en un punto de cada uno, se
multiplica por \(\Delta A\) y se suma sobre todos los rectángulos;
cuando el tamaño de la malla tiende a cero, ese límite (si existe) es
la integral doble
\cite[cap.~15]{stewart2018}\cite[cap.~13]{larson2015}\cite[cap.~11]{apostol}.
Cuando \(R\) es una región general (no necesariamente un rectángulo), el
\textbf{teorema de Fubini} permite calcular la integral doble como dos
integrales iteradas de una sola variable, integrando primero respecto
a una variable con la otra fija y luego respecto a la restante
\cite[pág.~982]{stewart2018}\cite[pág.~1210]{larson2015}\cite[pág.~348]{apostol}.\\

\noindent Teóricamente, existe una relación directa entre la intensidad acústica y la energía total en un espacio físico. La energía de un campo sonoro no es un valor concentrado en un punto aislado, sino que representa una magnitud que requiere sumar la densidad de energía \(\varepsilon\) a lo largo de un área o volumen específico utilizando herramientas de integración. La energía total \(E\) contenida en un espacio tridimensional se calcula integrando la densidad de energía sobre un volumen \(V\) mediante la fórmula general \cite{kinsler2000}:

\begin{equation}
E_{\text{total}} = \iiint_V \varepsilon \, dV.
\end{equation}

\noindent Puesto que en el contexto de este proyecto el análisis se concentra en un plano transversal (como el área \(A\) ocupada por el público o el escenario), la magnitud bidimensional que interesa integrar es la energía del campo sonoro proporcional a la presión eficaz al cuadrado, \(p_{\text{rms}}^2(x,y)\), definida en la sección~\ref{sec:campo_total}. La energía acústica total no se evalúa entonces en un solo punto, sino que se obtiene mediante una integral doble sobre dicha área \(A\):

\begin{equation}
E_{\text{total}}(A) = \iint_A p_{\text{rms}}^2(x,y)\, dA
= \iint_A E(x,y)\, dA.
\label{eq:energia_region}
\end{equation}

\noindent El objetivo específico correspondiente compara, mediante la
ecuación~\eqref{eq:energia_region}, la energía total obtenida
\emph{con} interferencia (integrando el campo real \(E(x,y)\) del
arreglo) frente a la que se obtendría sumando la energía que cada
fuente aportaría de manera independiente, $\sum_n \iint_R E_n(x,y)\,
dA$. La diferencia entre ambos resultados es, precisamente, la
manifestación integrada (sobre una región completa, y no solo en un
punto) del fenómeno de interferencia descrito en la
sección~\ref{sec:interferencia}: la energía no se conserva punto a
punto al sumar dos fuentes, sino que se redistribuye en el espacio.\\

\noindent Al igual que ocurre con los puntos críticos, la
ecuación~\eqref{eq:energia_region} no siempre admite una antiderivada
elemental para el integrando \(E(x,y)\), que combina términos \(1/r^2\)
con productos de funciones trigonométricas. Por ello, BONAIS aproxima
esta integral mediante \textbf{métodos numéricos de cuadratura}, en
particular la regla del punto medio, que consiste en subdividir la
región \(R\) en una malla fina de celdas rectangulares, evaluar \(E(x,y)\)
en el centro de cada una y sumar esos valores ponderados por el área
de la celda \cite[sec.~4.4]{burden2011} \cite[sec.~4.8]{burden2011}. Cuanto más fina es la malla, mejor
aproxima esta suma a la integral exacta de la
ecuación~\eqref{eq:energia_region}.


\noindent La figura~\ref{fig:integral_doble} ilustra este procedimiento sobre la zona del público con una malla de $8\times8$ celdas.

\begin{figure}[H]
\centering
\includegraphics[width=0.75\textwidth]{figuras/mt_integral_doble.png}
\caption{Regla del punto medio: la zona del público se divide en celdas de área $\Delta A$ y se evalúa la energía en el centro de cada una. BONAIS usa la misma idea con $48\times48$ celdas. Fuente: elaboración propia.}
\label{fig:integral_doble}
\end{figure}

\section{Antecedentes: arreglos direccionales de subwoofers}

\noindent Las secciones anteriores presentaron el sustento físico y matemático
de la interferencia entre fuentes puntuales, y cómo esa interferencia
se analiza con las herramientas del Cálculo III. Esta última sección
cierra el capítulo describiendo, como \emph{antecedente} del proyecto,
el estado actual de la técnica de ingeniería que ya aplica ese
principio en el sonido en vivo: los \emph{arreglos direccionales de
subwoofers}. A diferencia de los temas anteriores, este conocimiento
no proviene de un libro de texto, sino de la práctica documentada de
la industria del audio profesional, por lo que las fuentes consultadas
corresponden a publicaciones técnicas especializadas del gremio.\\

\noindent Como ya se estableció, un subwoofer aislado radia el sonido de manera
casi omnidireccional, por ser una fuente acústicamente compacta frente
a su longitud de onda. En un evento en vivo esto genera dos problemas
prácticos: el sonido grave se acumula en el escenario, donde puede
saturar los micrófonos y favorecer la retroalimentación, y la
cobertura de graves entre el público resulta desigual
\cite{fohonline_cardioid}. Como no es posible construir un subwoofer
físicamente direccional a esas frecuencias, la industria resuelve el
problema combinando varias unidades con una separación y un retardo
específicos entre ellas (exactamente los parámetros \(d\) y \(\tau\) de
la ecuación~\eqref{eq:fase_neta}), de modo que sus campos de presión
interfieran constructivamente hacia el público y destructivamente
hacia atrás \cite{eaw_cardioid}.\\

\noindent El principio detrás de esta técnica, conocida como arreglo de
\emph{gradiente} o, en su forma más simple, arreglo \emph{cardioide},
no es nuevo: su formulación se remonta al trabajo de Harry F. Olson
sobre micrófonos y altavoces de gradiente de presión, patentado a
comienzos de la década de 1970, en el que ya se establecía que
invertir la polaridad de una fuente y retardarla respecto a otra
produce un patrón de radiación direccional en bajas frecuencias
\cite{audioxpress_olson}. Lo que ha cambiado desde entonces es la
disponibilidad de procesadores digitales que permiten ajustar ese
retardo con precisión de fracciones de milisegundo, lo cual ha hecho
de los arreglos cardioides una práctica común en el sonido en vivo
contemporáneo \cite{fohonline_cardioid,eaw_cardioid}.\\

\noindent En la configuración más simple de dos subwoofers, la unidad trasera se
invierte en polaridad y se retarda un tiempo igual al que tarda el
sonido en recorrer la separación entre ambas, mientras que la
separación misma se fija, en la práctica, en una cuarta parte de la
longitud de onda de la frecuencia de interés \cite{eaw_cardioid}. Los
fabricantes reportan, para arreglos de este tipo, atenuaciones típicas
de hasta 15~dB hacia la parte trasera del arreglo respecto al frente
\cite{fohonline_cardioid}.

\noindent La figura~\ref{fig:polar} compara el patrón de radiación de dos subwoofers en fase con el de un arreglo cardioide de este tipo.

\begin{figure}[H]
\centering
\includegraphics[width=0.6\textwidth]{figuras/mt_polar.png}
\caption{Patrón polar a 80~Hz. Con dos subwoofers en fase el grave sale casi igual hacia todos lados; con el arreglo cardioide se concentra hacia el público y se anula hacia el escenario. Fuente: elaboración propia.}
\label{fig:polar}
\end{figure}

\noindent Este antecedente es, en últimas, el que
motiva la pregunta central del proyecto: determinar, a partir del
modelo matemático desarrollado en este capítulo, la separación y el
retardo que maximizan dicha atenuación para una frecuencia dada, y qué
tan sensible es ese resultado a los errores de ajuste que ocurren en
la práctica.